\scanpage{56}
\mspage{55}
2 - 4 - \hfill n° 392

$1\le i\le n$. Soit $x_0$ un élément de $K$ tel que $v_i(x_0)=q_i$ pour $1\le i\le n$, (chap. VI, § 7, n° 2, cor. 1 du th. 1). Alors $\underline a$ est l'ensemble des $x\in K$ tels que $v_i(xx_0^{-1})\ge0$ pour $1\le i\le n$. Ainsi $\underline a=Ax_0$.

Si $A$ est un anneau de Dedekind, on a vu, dans la démonstration du th. 1, que le groupe $D(A)$ des diviseurs de $A$ s'identifie au groupe $I(A)$ des idéaux fractionnaires $\underline a\ne(0)$, (comme $A$ est noethérien, tout idéal fractionnaire non nul est de type fini). Le groupe des \uline{classes de diviseurs} de $A$, (§ 1, n° 4), s'identifie alors au groupe des \uline{classes d'idéaux} de $A$, (chap. II, § 5, n° 7).

3. \uline{Décomposition des idéaux en produits d'idéaux premiers}.

Soient $A$ un anneau de Dedekind, $I(A)$ le groupe multiplicatif ordonné des idéaux fractionnaires non nuls de $A$, et $D(A)$ le groupe des diviseurs de $A$. L'isomorphisme $\underline a\longrightarrow\operatorname{div}\underline a$ de $I(A)$ sur $D(A)$ fait correspondre les diviseurs extrémaux aux idéaux premiers non nuls de $A$, (§ 1, n° 6, th. 3), donc le groupe multiplicatif $I(A)$ admet pour base l'ensemble des idéaux premiers non nuls de $A$, (§ 1, n° 3, th. 2). Autrement dit, \uline{tout idéal fractionnaire non nul} $\underline a$ \uline{de} $A$ \uline{admet une décomposition et une seule de la forme} :
\[
(1)\qquad\underline a=\prod\underline p^{n(\underline p)}
\]
où le produit est étendu aux idéaux premiers non nuls de $A$, les exposants $n(\underline p)$ étant nuls à l'exception d'un nombre fini d'entre eux. De plus $\underline a$ est entier si et seulement si les $n(\underline p)$ sont tous positifs. On dit que la relation (1) est la \uline{décomposition de} $\underline a$ \uline{en facteurs premiers}. En particulier, si $\underline a$ est un idéal principal $Ax$, on a, pour tout $\underline p$, $n(\underline p)=v_{\underline p}(x)$, où $v_{\underline p}$ désigne la valuation essentielle correspondant à $\underline p$ ; ceci résulte en effet de la formule (4) du § 1, n° 4. Soient
\[
\underline a=\prod\underline p^{m(\underline p)}\quad,\quad\underline b=\prod\underline p^{n(\underline p)}
\]
deux idéaux fractionnaires non nuls de $A$. On a alors :
